% Lesefassung zur Identifikation durch Multiplikationstafeln in Band 37.
% Aussagen und Nummern werden ausschließlich im Hauptband registriert.
\hypertarget{table-isomorphisms.reading}{}

\section*{Wann beschreiben zwei Tafeln dieselbe Halbgruppe?}

Eine endliche Halbgruppe lässt sich vollständig durch eine
Multiplikationstafel angeben. Die Zeile nennt den ersten Faktor, die
Spalte den zweiten; im Schnittpunkt steht das Produkt. Die Namen der
Elemente sind dabei frei wählbar. Zwei verschieden beschriftete Tafeln
können deshalb dieselbe algebraische Struktur beschreiben.

Gesucht ist eine Bijektion \(f\colon A\to B\), die jedes Produkt erhält:
\[
  f(x\star y)=f(x)\diamond f(y)\qquad(x,y\in A).
\]
Eine solche Abbildung heißt Isomorphismus. Die entscheidende Frage lautet
also: Lassen sich sämtliche Elementnamen so ersetzen, dass aus der einen
Multiplikation die andere wird?

Diese Lesefassung entwickelt daraus einen vollständigen endlichen Test.
Wir transportieren die Multiplikation zuerst auf einen nummerierten
Träger. Danach betrachten wir alle Umbenennungen und fassen ihre Tafeln
in einem Code zusammen. Gleiche Codes werden genau die isomorphen
Halbgruppen kennzeichnen. Die nummerierten Aussagen stehen weiterhin in
Band~37; ihre formalen Ableitungen sind in den
\TableIsomorphismsProofsLink{} zusammengestellt.

\section*{Von Elementnamen zu Zahlen}

Für \(n\) Elemente verwenden wir den Träger
\(S_n=\NatSeg n=\{0,\ldots,n-1\}\); bei \(n=0\) ist er leer.
Eine Nummerierung ist eine Bijektion
\[
  e\colon S_n\bij A.
\]
Zu ihr gehört die Zahlentafel
\[
  m_{\star,e}(i,j)=e^{-1}\bigl(e(i)\star e(j)\bigr).
\]
Man liest zuerst die beiden Elementnamen zu den Nummern \(i,j\),
bildet ihr Produkt in \(A\) und schreibt dessen Nummer in das Feld.
Die Abgeschlossenheit der Halbgruppe stellt sicher, dass das Produkt
wieder einen Namen in \(A\) besitzt. Die Bijektivität von \(e\) gibt
ihm genau eine Nummer. Dies ist
\FormulaRefAuto{FiniteSemigroupMultiplicationTableDef}[def].

Als durchgehendes Beispiel betrachten wir \(A=\{a,b,c\}\) mit der Tafel
\[
  \begin{array}{c|ccc}
    \star&a&b&c\\\hline
    a&b&c&b\\
    b&c&b&c\\
    c&b&c&b
  \end{array}.
\]
Diese Operation ist assoziativ. Um das ohne eine Liste von 27 Tripeln
zu sehen, setze \(q(a)=q(c)=1\) und \(q(b)=0\). Ein Produkt ist genau
dann \(b\), wenn die Summe der beiden \(q\)-Werte gerade ist;
andernfalls ist es \(c\). Insbesondere gilt
\[
  q(x\star y)\equiv q(x)+q(y)\pmod 2.
\]
Beide Klammerungen eines Dreifachprodukts liegen in \(\{b,c\}\), und
beide besitzen den \(q\)-Wert der Summe
\(q(x)+q(y)+q(z)\) modulo \(2\). Auf \(\{b,c\}\) unterscheidet
\(q\) die beiden Elemente. Daher sind die beiden Klammerungen gleich.

Mit \(e(0)=a\), \(e(1)=b\), \(e(2)=c\) erhalten wir
\[
  \begin{array}{c|ccc}
    m&0&1&2\\\hline
    0&1&2&1\\
    1&2&1&2\\
    2&1&2&1
  \end{array}.
\]
Hier und im Folgenden kürzen wir \(m_{\star,e}\) zu \(m\) ab.
Beispielsweise steht im Feld \((0,1)\) die Zahl \(2\), weil
\(a\star b=c=e(2)\). Die Zahl \(0\) ist lediglich der Name von
\(a\); sie bezeichnet kein Nullelement.

Die Nummerierung erhält die gesamte Multiplikation:
\[
  e(m(i,j))=e(i)\star e(j).
\]
Damit überträgt sie auch die Assoziativität. Tatsächlich werden
\(m(m(i,j),k)\) und \(m(i,m(j,k))\) unter \(e\) auf die beiden
gleichen Klammerungen in \(A\) abgebildet; aus der Injektivität von
\(e\) folgt ihre Gleichheit. Die Zahlentafel ist folglich wieder eine
endliche Halbgruppe, und \(e\) ist ein Isomorphismus. Dies beweist den
Inhalt von
\FormulaRefAuto{FiniteSemigroupMultiplicationTableTransport}[thm].
\TableIsomorphismsProofLink{FiniteSemigroupMultiplicationTableTransport}.

\section*{Eine Umbenennung betrifft auch die Ergebnisse}

Eine Permutation \(\pi\colon S_n\bij S_n\) gibt jedem bisherigen
Index einen neuen Namen. Wird aus \(i\) der Name \(\pi(i)\), dann
muss auch das Produkt \(m(i,j)\) seinen neuen Namen erhalten. Für die
umbenannte Tafel \(d\) gilt deshalb
\[
  d(\pi(i),\pi(j))=\pi(m(i,j))\qquad(i,j\in S_n).
  \tag{U}\label{table-isomorphisms-reading-relabel}
\]
Dieselbe Permutation wirkt auf die Zeilen, die Spalten und die
Tafeleinträge. Werden nur die Zeilen und Spalten verschoben, behalten
die Ergebnisse ihre alten Namen; damit ist die Umbenennung im
Allgemeinen noch nicht vollzogen.

Im Beispiel vertauschen wir \(0\) und \(1\) und lassen \(2\) fest.
Die neue Tafel lautet
\[
  \begin{array}{c|ccc}
    d&0&1&2\\\hline
    0&0&2&2\\
    1&2&0&0\\
    2&2&0&0
  \end{array}.
\]
So wird \(m(0,0)=1\) zu \(d(1,1)=0\), und aus
\(m(0,1)=2\) wird \(d(1,0)=2\). Das idempotente Element \(b\)
hieß zuvor \(1\) und heißt jetzt \(0\). Auch diese Eigenschaft
wurde mittransportiert.

Um ein beliebiges Feld \((u,v)\) der neuen Tafel auszurechnen,
gehen wir zunächst zu den alten Namen zurück. Damit wird
\eqref{table-isomorphisms-reading-relabel} zur Formel
\[
  m^\pi(u,v)
   =\pi\bigl(m(\pi^{-1}(u),\pi^{-1}(v))\bigr).
\]
Sie definiert die umbenannte Operation in
\FormulaRefAuto{FiniteSemigroupRelabelledOperationDef}[def].
Die Schreibweisen \(d=m^\pi\) und
\eqref{table-isomorphisms-reading-relabel} sagen also dasselbe.

Die Produktgleichungen sind zugleich ein vollständiger Test für eine
\emph{vorgegebene} Bijektion: Erfüllt \(\pi\) sie in allen Feldern,
so ist \(\pi\) ein Isomorphismus; jeder Isomorphismus muss sie erfüllen.
Mehr verlangt die Definition eines Isomorphismus hier nicht.
Das ist
\FormulaRefAuto{FiniteSemigroupTableRelabelCriterion}[thm].
\TableIsomorphismsProofLink{FiniteSemigroupTableRelabelCriterion}.

Für zwei beliebige Träger gleicher Größe wählen wir Nummerierungen
\(e_A\colon S_n\bij A\) und \(e_B\colon S_n\bij B\). Ein
Isomorphismus \(f\colon A\to B\) erscheint in den Tafeln als
\[
  \pi=e_B^{-1}\circ f\circ e_A.
\]
Umgekehrt ergibt eine passende Umbenennung der Tafeln den Isomorphismus
\[
  f=e_B\circ\pi\circ e_A^{-1}.
\]
Damit reicht die Suche nach Permutationen des festen Zahlenträgers
für den Vergleich beliebiger endlicher Halbgruppen aus.

\section*{Orbits und die Symmetrien einer Tafel}

Auf einem festen Träger \(S\) sammeln wir die assoziativen binären
Operationen in \(\operatorname{Assoc}(S)\) und die Permutationen in
\(\operatorname{Sym}(S)\). Dies sind die Mengen aus
\FormulaRefAuto{FiniteSemigroupAssociativeOperationsDef}[def] und
\FormulaRefAuto{FiniteSemigroupSymmetricSetDef}[def]. Die Umbenennung
ist mit dem Hintereinanderausführen von Permutationen verträglich:
\[
  m^{\Id_S}=m,\qquad
  (m^\pi)^\sigma=m^{\sigma\circ\pi}.
\]
Zuerst nach \(\pi\), dann nach \(\sigma\) umzubenennen hat also
denselben Effekt wie die gemeinsame Umbenennung \(\sigma\circ\pi\).
Jede dieser Tafeln bleibt assoziativ. Diese drei Aussagen bilden den
Kern von
\FormulaRefAuto{FiniteSemigroupRelabellingAction}[thm].
\TableIsomorphismsProofLink{FiniteSemigroupRelabellingAction}.

Der Hauptband schreibt
\(\operatorname{Relab}_S(\pi;m,d)\), wenn \(\pi\) eine solche
Umbenennung von \(m\) nach \(d\) ist. Diese Relation speichert
zugleich die Permutation und die Produktgleichungen. Dass sie genau
\(d=m^\pi\) ausdrückt, zeigt
\FormulaRefAuto{FiniteSemigroupRelabelGraphCharacterization}[thm].
\TableIsomorphismsProofLink{FiniteSemigroupRelabelGraphCharacterization}.

Der \emph{Umbenennungsorbit} einer Tafel ist die Menge aller Tafeln,
die auf diese Weise entstehen:
\[
  \operatorname{Orb}_S(m)
   =\{m^\pi\mid\pi\in\operatorname{Sym}(S)\}.
\]
Dies ist die Beschreibung von
\FormulaRefAuto{FiniteSemigroupRelabelOrbitDef}[def] mithilfe der
Umbenennungsformel. Auf einem endlichen festen Träger sind zwei
Operationen genau dann isomorph, wenn sie im selben Orbit liegen.
Die Hinrichtung benutzt die Produktgleichungen des Isomorphismus,
die Rückrichtung die Permutation, mit der umbenannt wurde. Somit
sind die Orbits genau die Isomorphieklassen:
\FormulaRefAuto{FiniteSemigroupRelabelOrbitsAreIsomorphismClasses}[thm].
\TableIsomorphismsProofLink{FiniteSemigroupRelabelOrbitsAreIsomorphismClasses}.

Eine Permutation kann die Tafel auch unverändert lassen. Dann heißt
sie \emph{Automorphismus}. Gemäß
\FormulaRefAuto{FiniteSemigroupAutomorphismGroupDef}[def] gilt
\[
  \operatorname{Aut}(S,m)
   =\{\pi\in\operatorname{Sym}(S)\mid m^\pi=m\}.
\]
Diese Permutationen bilden unter Komposition eine Gruppe: Die
Identität lässt die Tafel fest, zwei solche Umbenennungen tun dies
auch hintereinander, und das Rückgängigmachen einer solchen
Umbenennung lässt die Tafel ebenfalls fest. Alle drei Eigenschaften
folgen unmittelbar aus dem Wirkungsgesetz.

Unser Dreierbeispiel besitzt nur den identischen Automorphismus.
Das Element \(b\) ist das einzige idempotente Element und muss daher
fest bleiben. Die Menge der Produktwerte ist \(\{b,c\}\) und wird
von jedem Automorphismus auf sich abgebildet. Somit bleibt auch
\(c\) fest und schließlich \(a\). Deshalb liefern die sechs
Permutationen hier sechs verschiedene Tafeln: Wären zwei Ergebnisse
gleich, so wäre die erste Umbenennung gefolgt von der Umkehrung der
zweiten ein nichttrivialer Automorphismus.

Anders ist es bei der Nullhalbgruppe auf \(\{a,b,c\}\), in der
\emph{jedes} Produkt \(b\) ist. Hier darf man \(a\) und \(c\)
vertauschen. Genau die beiden Permutationen, die \(b\) festlassen,
sind Automorphismen. Unter den sechs Umbenennungen entstehen nur
drei verschiedene Tafeln, je nachdem, welchen Namen der konstante
Produktwert erhält. Die Zahl der Umbenennungen und die Zahl ihrer
verschiedenen Ergebnisse können also auseinanderfallen.

\section*{Eine Tafel als Tupel lesen}

Wir lesen die Einträge zeilenweise und lassen die Zeilen- und
Spaltenköpfe weg. Aus einer Tafel auf \(S_n\) entsteht so ein Tupel
\(w_n(m)\) der Länge \(n^2\), wie in
\FormulaRefAuto{FiniteSemigroupTableTupleDef}[def]. Für die beiden
Tafeln unseres Beispiels gilt
\[
\begin{aligned}
  w_3(m)&=(1,2,1,\ 2,1,2,\ 1,2,1),\\
  w_3(d)&=(0,2,2,\ 2,0,0,\ 2,0,0).
\end{aligned}
\]
Die Leerstellen gliedern hier nur die drei Zeilen. Es handelt sich
um Tupel von Zahlen; auch mehrstellige Einträge bleiben durch die
Trennzeichen eindeutig.

Gleiche Tupel geben an jeder Position denselben Eintrag an.
Folglich bestimmen sie dieselbe nummerierte Operation. Das ist
\FormulaRefAuto{FiniteSemigroupTableTupleInjective}[thm].
\TableIsomorphismsProofLink{FiniteSemigroupTableTupleInjective}.
Verschiedene Tupel können dagegen isomorphe Halbgruppen darstellen,
wie \(m\) und \(d\) zeigen. Ein einzelnes, beliebig nummeriertes
Tupel genügt daher noch nicht als nummerierungsunabhängiger Code.

\section*{Der vollständige kanonische Code}

Wir sammeln stattdessen sämtliche Tafeltupel, die durch
Nummerierungen derselben Halbgruppe entstehen. Zusammen mit der
Elementzahl ergibt dies den Code
\[
\begin{aligned}
  \operatorname{can}(A,\star)
   &=\bigl(n,\mathcal W(A,\star)\bigr),
   \qquad n=\FiniteCard(A),\\
  \mathcal W(A,\star)
   &=\{w_n(m_{\star,e})\mid e\colon S_n\bij A\}.
\end{aligned}
\]
Hier steht \(\mathcal W(A,\star)\) als Abkürzung für die zweite
Komponente von
\FormulaRefAuto{FiniteSemigroupCanonicalTableCodeDef}[def]. Die
Menge ist endlich und nichtleer. Das gilt auch für den leeren
Träger: Er besitzt genau eine Nummerierung und genau ein leeres
Tafeltupel.

Im Code wird kein bestimmter Ausgangsname bevorzugt. Beginnen wir
mit einer anderen Nummerierung, erhalten wir dieselbe Menge aller
Möglichkeiten. Umgekehrt enthält der Code genügend Information,
um die Halbgruppe bis auf Isomorphie zurückzugewinnen:
\[
  (A,\star)\cong(B,\diamond)
  \quad\Longleftrightarrow\quad
  \operatorname{can}(A,\star)=\operatorname{can}(B,\diamond).
\]
Dies ist
\FormulaRefAuto{FiniteSemigroupCanonicalTableCodeComplete}[thm].
Beide Richtungen lassen sich direkt verstehen.

Ist \(f\colon A\to B\) ein Isomorphismus, so macht jede Nummerierung
\(e_A\) von \(A\) die Abbildung \(e_B=f\circ e_A\) zu einer
Nummerierung von \(B\). Unter diesen Nummerierungen sind alle
Tafeleinträge gleich, denn \(f\) erhält die Produkte. Jedes Tupel
von \(A\) kommt deshalb bei \(B\) vor. Mit \(f^{-1}\) folgt auch
die umgekehrte Inklusion. Die Elementzahlen stimmen ohnehin überein,
also sind die Codes gleich.

Sind umgekehrt die Codes gleich, haben \(A\) und \(B\) dieselbe
Elementzahl. Wähle eine Nummerierung \(e_A\). Ihr Tafeltupel kommt
auch bei einer Nummerierung \(e_B\) von \(B\) vor. Die Tupelgleichheit
liefert die Gleichheit der nummerierten Operationen. Für
\(x=e_A(i)\) und \(y=e_A(j)\) folgt daher mit
\(f=e_B\circ e_A^{-1}\)
\[
\begin{aligned}
  f(x\star y)
   &=e_B(m_{\star,e_A}(i,j))\\
   &=e_B(m_{\diamond,e_B}(i,j))
     =f(x)\diamond f(y).
\end{aligned}
\]
Die Abbildung \(f\) ist bijektiv und erhält jedes Produkt. Damit ist
der gesuchte Isomorphismus konstruiert.
\TableIsomorphismsProofLink{FiniteSemigroupCanonicalTableCodeComplete}.

\section*{Für die Rechnung genügt das kleinste Tupel}

Die ganze Tupelmenge lässt sich aus einer einzigen Ausgangstafel
berechnen. Jede weitere Nummerierung hat die Form \(e\circ\rho\),
wobei \(\rho\) eine Permutation von \(S_n\) ist. Dabei gilt
\[
  m_{\star,e\circ\rho}=m_{\star,e}^{\rho^{-1}}.
\]
Die Umkehrung im Exponenten ist wesentlich: Die neue Nummerierung
weist einer Nummer ein Element zu; die Umbenennung weist einer alten
Nummer ihren neuen Namen zu. Da auch \(\rho^{-1}\) alle Permutationen
durchläuft, erhalten wir genau die Tafeltupel des Umbenennungsorbits.

Man kann diese endliche Menge durch ihr lexikographisch kleinstes
Tupel vertreten. Beim Vergleich zweier Tupel entscheidet die erste
Position, an der ihre Einträge verschieden sind. Die kleinere Zahl
an dieser Position ergibt das kleinere Tupel. Wir speichern die
Elementzahl zusammen mit diesem Minimum. Diese rechnerische Form
ist vom formal definierten Code mit seiner ganzen Tupelmenge zu
unterscheiden; sie enthält dieselbe Isomorphieinformation.

Warum genügt das Minimum? Haben zwei Orbits dasselbe kleinste Tupel,
enthalten sie dieselbe nummerierte Tafel. Von dieser gemeinsamen
Tafel aus erreicht man durch Umbenennen jede Tafel beider Orbits.
Die Orbits sind deshalb gleich. Die umgekehrte Richtung ist
unmittelbar: Gleiche Orbits haben dasselbe Minimum.

Im Dreierbeispiel ergibt die neue Nummerierung
\[
  e'(0)=b,\qquad e'(1)=c,\qquad e'(2)=a
\]
die Tafel
\[
  \begin{array}{c|ccc}
    m_{\star,e'}&0&1&2\\\hline
    0&0&1&1\\
    1&1&0&0\\
    2&1&0&0
  \end{array}.
\]
Ihr Tupel
\[
  (0,1,1,\ 1,0,0,\ 1,0,0)
\]
ist das kleinste der sechs Kandidaten. Nur zwei Kandidaten beginnen
mit \(0\): Dazu muss das einzige idempotente Element \(b\) die
Nummer \(0\) erhalten. Erhält anschließend \(c\) die Nummer \(1\),
ist der zweite Eintrag \(1\); erhält \(a\) diese Nummer, ist der
zweite Eintrag \(2\). Damit ist die Wahl des Minimums bereits
entschieden.

Ein direkter Algorithmus für eine vorliegende Halbgruppentafel ist
nun vollständig beschrieben:
\begin{enumerate}
  \item Bestimme die Trägergröße \(n\) und wähle eine Nummerierung.
  \item Durchlaufe alle Permutationen \(\pi\) von \(S_n\).
  \item Berechne jeweils \(m^\pi\) und lies sein Tafeltupel.
  \item Behalte das kleinste bisher gefundene Tupel sowie eine
    Permutation, die es erzeugt.
  \item Vergleiche bei zwei Halbgruppen die Paare aus Trägergröße
    und kleinstem Tupel.
\end{enumerate}
Unterschiedliche Paare schließen Isomorphie aus; gleiche Paare
beweisen sie. Voraussetzung ist, dass die eingegebenen Tafeln
binäre assoziative Operationen darstellen. Bei ungeprüften Daten
sind diese Bedingungen vorab zu prüfen.

Die gespeicherten Permutationen liefern außerdem einen konkreten
Isomorphismus. Haben die Ausgangstafeln \(m_A,m_B\) unter
\(\pi_A,\pi_B\) dieselbe minimale Tafel, so bildet
\(\pi_B^{-1}\circ\pi_A\) die erste Ausgangstafel isomorph auf die
zweite ab. Auf den ursprünglichen Trägern ist dies
\[
  f=e_B\circ\pi_B^{-1}\circ\pi_A\circ e_A^{-1}.
\]
Der Algorithmus kann also neben dem Ja oder Nein auch einen Zeugen
für die Isomorphie ausgeben.

Es gibt \(n!\) Permutationen und je \(n^2\) Tafelfelder. Die direkte
Berechnung benötigt somit eine Anzahl elementarer Tafelarbeiten
von der Größenordnung \(n!\,n^2\). Sie endet für jedes feste \(n\),
kann aber schon bei mäßig großen Trägern aufwendig werden.
Automorphismen können dazu führen, dass verschiedene Permutationen
wiederholt dieselbe Tafel liefern.

\section*{Was schnelle Merkmale leisten}

Vor der vollständigen Suche lohnt sich der Vergleich leicht
berechenbarer Eigenschaften. Beispielsweise muss ein Isomorphismus
idempotente Elemente auf idempotente Elemente abbilden:
\[
  x^2=x\quad\Longrightarrow\quad f(x)^2=f(x^2)=f(x).
\]
Mit dem inversen Isomorphismus folgt die Rückrichtung. Deshalb
stimmen die Anzahlen der Idempotenten überein. Auch Kommutativität
und die Existenz eines Nullelements werden erhalten. Ein Unterschied
in einem solchen Merkmal beendet die Suche sofort.

Gleiche Merkmale allein reichen dagegen nicht. Die beiden Tafeln
\[
  \begin{array}{c|cc}
    m_L&0&1\\\hline
    0&0&0\\
    1&1&1
  \end{array}
  \qquad
  \begin{array}{c|cc}
    m_R&0&1\\\hline
    0&0&1\\
    1&0&1
  \end{array}
\]
erfüllen die Regeln \(m_L(x,y)=x\) und \(m_R(x,y)=y\).
Beide Operationen sind assoziativ, besitzen zwei idempotente
Elemente und sind nicht kommutativ. Dennoch sind sie nicht
isomorph. Ein Isomorphismus \(f\) würde für verschiedene \(x,y\)
liefern
\[
  f(x)=f(m_L(x,y))=m_R(f(x),f(y))=f(y),
\]
im Widerspruch zu seiner Injektivität. Ihre kleinsten Tafeltupel
sind entsprechend verschieden: \((0,0,1,1)\) und \((0,1,0,1)\).
Die einfachen Merkmale sind nützliche Vorfilter; den vollständigen
Nachweis liefern eine passende Umbenennung oder die gleichen Codes.

\section*{Die Grenze beim Rückweg aus der Potenzhalbgruppe}

Das Tafelverfahren entscheidet den Vergleich der Halbgruppen,
deren Tafeln eingegeben werden. Man kann es daher auch auf
Potenzhalbgruppen anwenden. In einer Potenzhalbgruppe sind die
Elemente jedoch selbst nichtleere Teilmengen, und ihre Multiplikation
lautet
\[
  XY=\{xy\mid x\in X,\ y\in Y\}.
\]
Ein abstrakter Isomorphismus dieser Halbgruppen muss die
Multiplikation erhalten. Die Größen der Teilmengen oder ihr
Einermengenstatus sind damit noch nicht als feste Markierungen
gegeben.

Das zeigt bereits die Nullhalbgruppe \(S=\{0,a\}\), in der jedes
Produkt \(0\) ist. Ihre Potenzhalbgruppe besitzt die drei Elemente
\[
  \{0\},\qquad\{a\},\qquad\{0,a\},
\]
und jedes ihrer Produkte ist \(\{0\}\). Die Bijektion, die
\(\{a\}\) und \(\{0,a\}\) vertauscht und \(\{0\}\) festlässt,
ist deshalb ein Automorphismus. Sie vertauscht eine Einermenge mit
einer Zweiermenge.

Wer Potenzhalbgruppen als abstrakte Halbgruppen vergleicht, muss
solche Umbenennungen zulassen. Eine zusätzliche Markierung aller
Einermengen würde einen anderen Test ergeben. Aus gleichen
kanonischen Codes der Potenzhalbgruppen folgt zunächst deren
Isomorphie. Wie daraus ein Isomorphismus der Ausgangshalbgruppen
gewonnen werden kann, ist eine weitere Rekonstruktionsfrage. Der
Tafelcode allein enthält dafür noch keinen allgemeinen Beweis.

Die hier behandelte Aufgabe ist damit präzise gelöst: Bei gegebenen
endlichen Halbgruppentafeln lässt sich Isomorphie durch vollständiges
Umbenennen entscheiden. Die Rekonstruktion einer Grundhalbgruppe
aus ihrer Potenzhalbgruppe wird im anschließenden Kapitel von
Band~37 gesondert behandelt.

\medskip
\noindent\TableIsomorphismsMainLink{}
\enspace\textperiodcentered\enspace
\TableIsomorphismsProofsLink{}
